\documentclass[9pt]{extarticle}
\usepackage[a4paper,landscape,margin=1cm]{geometry}
\usepackage{multicol}
\usepackage{xcolor}
\usepackage{titlesec}
\usepackage{enumitem}
\usepackage{tabularx}
\usepackage{booktabs}
\usepackage{fancyhdr}
\usepackage{needspace}
\usepackage[T1]{fontenc}
\usepackage{inconsolata}
\usepackage{helvet}
\renewcommand{\familydefault}{\sfdefault}

% ── palette ────────────────────────────────────────────────
\definecolor{ink}{HTML}{26221C}
\definecolor{accent}{HTML}{7A2E2E}
\definecolor{faint}{HTML}{6F675C}
\definecolor{rule}{HTML}{D9D2C5}
\definecolor{codebg}{HTML}{F3EFE8}
\color{ink}

% ── section styling ────────────────────────────────────────
\titleformat{\section}{\normalfont\bfseries\color{accent}\small}{}{0pt}{\MakeUppercase}
\titlespacing{\section}{0pt}{6pt}{2pt}
\titleformat{\subsection}{\normalfont\bfseries\footnotesize}{}{0pt}{}
\titlespacing{\subsection}{0pt}{4pt}{1pt}

\setlist[description]{leftmargin=0pt,itemsep=1pt,parsep=0pt,topsep=1pt,font=\normalfont}
\setlength{\parindent}{0pt}
\setlength{\columnsep}{1.2em}
\setlength{\columnseprule}{0.4pt}
\renewcommand{\columnseprulecolor}{\color{rule}}

\pagestyle{fancy}
\fancyhf{}
\renewcommand{\headrulewidth}{0pt}
\renewcommand{\footrulewidth}{0.4pt}
\renewcommand{\footrule}{\vskip1pt{\color{rule}\hrule height 0.4pt}\vskip2pt}
\fancyfoot[L]{\footnotesize\color{faint}\textbf{lit} — the Literium CLI · repository format~1 · segmentation \texttt{uax29-16.0.0}}
\fancyfoot[R]{\footnotesize\color{faint}\thepage}

% ── code + command helpers ─────────────────────────────────
\newcommand{\cmd}[1]{\colorbox{codebg}{\texttt{\small #1}}}
\newcommand{\ci}[1]{\texttt{#1}}                       % inline code
\newcommand{\opt}[1]{\textcolor{faint}{\texttt{#1}}}   % optional/args
% a command entry: name + forms (left), gloss (right)
\newcommand{\entry}[2]{\needspace{2\baselineskip}\item[\ci{\textbf{#1}}]#2}

\begin{document}
\raggedright

{\LARGE\bfseries\color{accent} lit}\quad{\large command reference}%
\hfill{\footnotesize\color{faint} version control for literary texts — the atom is the sentence (prose) or the line (poetry), not the line of code}

\vspace{2pt}{\color{rule}\hrule height 0.6pt}\vspace{4pt}

\begin{multicols}{3}

% ═══════════════════════════════════════════════════════════
\section{The model in one breath}
A repository lives in \ci{.lit/}. Text is the truth; everything
else derives deterministically under the pinned segmentation.
Six content-addressed object types (SHA-256): \textbf{atom}
(a sentence/line), \textbf{fasciculus} (a paragraph/stanza — a
list of atom hashes), \textbf{spina} (the document skeleton),
\textbf{atom-spina}, \textbf{volumen} (a document's bytes), and
\textbf{version} (identity trios + a mutable manifest). A
\textbf{history} is an append-only log of \textbf{inscriptiones}
(version, timestamp, message). The one \textbf{historia} is
canonical; \textbf{experiments} are alternate lines. Objects
live in the \textbf{theca} (\ci{.lit/theca/}).

The everyday loop: edit files $\rightarrow$ \ci{conscribe} onto
the \textbf{tabula} (scratchpad) $\rightarrow$ \ci{inscribe}
onto a history.

% ═══════════════════════════════════════════════════════════
\section{Recording work}
\begin{description}
\entry{conscribe}{Snapshot the working tree onto the tabula.
Creates the repository on first use.\\
\opt{--mode prose|poetry}\; segmentation mode at creation
(default \ci{prose}).\quad\opt{--yes}\; skip prompts.}
\entry{inscribe}{Finalize the tabula (or a named version) onto
a history — appends an inscriptio.\\
\opt{<target>}\; history to append to (default \ci{historia}).
Must already exist.\\
\opt{-<n>}\; take the n-th pending revision;\;
\opt{--all}\; take them all.\\
\opt{-m,\;--message <msg>}\; the inscriptio message.\\
\opt{-f,\;--fons,\;--from <src>}\; inscribe an existing version
(adoption / re-inscribe) instead of the tabula.}
\entry{status}{List pending revisions relative to the latest
historia version.\quad\opt{<target-history>}.}
\entry{exscribe}{Materialize a version into the working tree
(document-set fidelity: files absent from the version are
deleted).\\
\ci{exscribe <src> [<path>]}\; whole version, or one document.\quad
\opt{--force}\; overwrite unrecorded changes.}
\entry{tabula}{Inspect the tabula.\quad
\ci{tabula discard}\; empty it.}
\end{description}

% ═══════════════════════════════════════════════════════════
\section{Histories \& experiments}
\begin{description}
\entry{experiment}{Manage alternate histories alongside
historia.\\
\ci{experiment create <name>}\;·\;\ci{experiment delete <name>}\;·\;\ci{experiment list}.\\
Names: not \ci{historia}, not 64-hex, no \ci{/}\,·\,whitespace\,·\,control, no leading \ci{.}}
\entry{contexe}{Combining merge — weave a source history/version
into the working tree, then \ci{conscribe}+\ci{inscribe}.\\
\ci{contexe <source> [<target>]}\; (target default \ci{historia}).\\
\opt{--base <version>}\; the three-way base.\quad
\opt{--force}.}
\entry{conflicts}{List entries in the contextio ledger
(\ci{.lit/contextio}).\quad\ci{conflicts clear}\; discard it.\\
\emph{Conflicts materialize inline as plain text — no markers;
resolve by ordinary editing, then conscribe.}}
\end{description}
\emph{Adoption (fast-forward) is \ci{inscribe --fons <version>};
there is no cherry-pick / rebase / revert — \ci{--fons} an
earlier version re-inscribes it.}

% ═══════════════════════════════════════════════════════════
\section{Inspection}
\begin{description}
\entry{show}{View a history or a version.\\
\ci{show [<history>|<version> [<path> [--atoms|<coord>]]]}\\
\opt{--atoms}\; list a document's atoms;\;
\opt{--hash}\; show hashes;\;
\opt{<coord>}\; a \ci{gNaM} coordinate (grex \emph{N}, atom
\emph{M}, e.g.\ \ci{g1a2}).}
\entry{diff}{Compare two versions, or the tabula / working tree
against the latest.\\
\ci{diff [<version> [<version>]]}.}
\entry{stemma}{A version's life-path across the histories; with
a coordinate, an atom's first/last inscription.\\
\ci{stemma [<version> [<path> <coord>]]}.}
\entry{fsck}{Verify repository integrity (the object walk).}
\entry{gc}{Delete unreachable objects.\quad\opt{--dry-run}.}
\end{description}

% ═══════════════════════════════════════════════════════════
\section{Correspondence}
\emph{Collaboration is mail: share a version's raw files, edit,
mail back, diff, inscribe. No account needed.}
\begin{description}
\entry{exemplar}{Pack a version's documents (+ an \ci{.exemplar}
metadata file: marker, basis, segmentation) into a
\ci{.tgz}.\\
\ci{exemplar [<source>] [-o <file>]}\; — with no source, in an
unpacked exemplar dir, repacks the return trip.}
\entry{collate}{Ingest a returned exemplar into the working
tree; reports the \emph{basis} for review.\quad
\ci{collate <file> [--force]}.\\
\emph{Then: \ci{diff <basis>}, \ci{conscribe}+\ci{inscribe} —
or \ci{contexe historia --base <basis>} if historia moved.}}
\end{description}

% ═══════════════════════════════════════════════════════════
\section{Litorium account}
\emph{Passwordless: an email link, first sign-in creates the
account.}
\begin{description}
\entry{account login}{\ci{account login [--email <addr>]} — a
one-time link arrives by mail; compare the code, open it. Then
prompts for a handle.}
\entry{account handle}{\ci{account handle [<name>]} — show or
claim the account handle (set-once; lowercase alnum + single
hyphens, min 12 chars self-serve).}
\entry{account status}{Signed-in email + handle, or ``not
signed in''.}
\entry{account logout}{Revoke the refresh token; erase the
credential file.}
\end{description}

% ═══════════════════════════════════════════════════════════
\section{Sync (litorium)}
\begin{description}
\entry{push}{Sync the historia to litorium; the first push
creates the project.\quad\opt{--remote <url>} (first push
only).}
\entry{pull}{Fast-forward the historia from litorium (divergence
is an error — push first).}
\entry{clone}{Fetch a project into a new directory.\\
\ci{clone <project-id> [<dir>] [--remote <url>]}.}
\end{description}

% ═══════════════════════════════════════════════════════════
\section{Publish \& share}
\begin{description}
\entry{publish}{Contribute the historia's (or one version's)
atoms to the public corpus — permanent, consent-ful.\\
\ci{publish [<version>] [--yes]}.}
\entry{share}{\ci{share} — export the historia to stdout.\\
\ci{share <version> [--expires <days>]} — mint a revocable
\ci{ref.litorium.org} link (default 30 d, max 365).}
\entry{shares}{List active ephemeral shares.}
\entry{opus}{\ci{opus <slug> <version>} — register a permanent
\ci{opus.litorium.org/<handle>/<slug>} permalink (needs a
handle; write-once).}
\entry{opera}{List registered opera.}
\entry{retract}{\ci{retract --link <token>} — revoke a share.\\
\ci{retract <version>} — drop an inscriptio (locally, and the
remote historia by compare-and-swap).}
\end{description}

% ═══════════════════════════════════════════════════════════
\section{Source \& version resolution}
Wherever a command takes a \ci{<source>} or \ci{<version>}:
\begin{description}[leftmargin=1em,labelsep=0.4em]
\item[$\cdot$] 64 hex chars $\rightarrow$ a version hash (must
exist).
\item[$\cdot$] 6–63 hex chars $\rightarrow$ a unique version
prefix (unless it names a history).
\item[$\cdot$] \ci{historia} or an experiment name
$\rightarrow$ that history's latest inscribed version.
\end{description}

% ═══════════════════════════════════════════════════════════
\section{git peers}
Every verb has an equal git-name peer (type either):
\begin{description}[leftmargin=1em,labelsep=0.4em,font=\normalfont]
\item[$\cdot$] \ci{add}, \ci{stage}$\to$conscribe\quad
\ci{commit}$\to$inscribe
\item[$\cdot$] \ci{checkout}, \ci{restore}$\to$exscribe\quad
\ci{merge}$\to$contexe
\item[$\cdot$] \ci{branch}$\to$experiment\quad
\ci{log}$\to$show\quad\ci{blame}$\to$stemma
\item[$\cdot$] same word both worlds: \ci{status}, \ci{diff},
\ci{gc}, \ci{fsck}, \ci{push}, \ci{pull}, \ci{clone}
\end{description}
Omitted git verbs (\ci{init}, \ci{revert}, \ci{cherry-pick},
\ci{rebase}, \ci{switch}, \ci{stash}, \ci{fetch}, \ci{tag},
\ci{remote}) print a one-line redirect.

\section{Aliases (Latin / Greek)}
The CLI accepts classical forms:
\begin{description}[leftmargin=1em,labelsep=0.4em,font=\normalfont]
\item[$\cdot$] \ci{conscribo}$\to$conscribe\quad
\ci{inscribo}, \ci{epigrapho}$\to$inscribe
\item[$\cdot$] \ci{exscribo}, \ci{ekgrapho}$\to$exscribe
\item[$\cdot$] \ci{contexo}, \ci{synhyphaino}$\to$contexe
\item[$\cdot$] \ci{experimentum}$\to$experiment\quad
\ci{confero}$\to$collate
\end{description}

% ═══════════════════════════════════════════════════════════
\section{Files \& layout}
\ci{.lit/} contains exactly:
\begin{description}[leftmargin=1em,labelsep=0.4em,font=\normalfont]
\item[$\cdot$] \ci{format}\; ``1'';\quad
\ci{segmentation}\; the pinned version.
\item[$\cdot$] \ci{historia}\; the canonical history file.
\item[$\cdot$] \ci{experimenta/<name>}\; one file per
experiment.
\item[$\cdot$] \ci{tabula}\; the scratchpad (absent = empty).
\item[$\cdot$] \ci{contextio}\; the merge conflict ledger
(absent = none).
\item[$\cdot$] \ci{theca/}\; the object store.
\item[$\cdot$] \ci{litorium}\; the remote record
\ci{\{remote, project, watermark\}} (synced repos).
\end{description}
Attributes: \ci{.litattributes} sets per-path mode;
\ci{.litignore} (gitignore-style) excludes paths.\\
Credentials: \ci{\$XDG\_CONFIG\_HOME/lit/credentials} (default
\ci{\textasciitilde/.config/lit/credentials}, mode 0600):
\ci{issuer}, \ci{email}, \ci{refresh\_token}, \ci{handle}.

% ═══════════════════════════════════════════════════════════
\section{Environment}
\begin{description}[leftmargin=1em,labelsep=0.4em,font=\normalfont]
\item[$\cdot$] \ci{LIT\_ID\_URL}\; identity host (default
\ci{id.litorium.org}).
\item[$\cdot$] \ci{LIT\_GATE\_URL}\; storage gate (default
\ci{p.litorium.org}).
\item[$\cdot$] \ci{LIT\_OPERA\_URL}\; sharing service (default
\ci{ref.litorium.org}).
\item[$\cdot$] \ci{LIT\_TOKEN}\; a bearer token overriding the
signed-in account (CI / static tokens).
\item[$\cdot$] \ci{XDG\_CONFIG\_HOME}\; credential-file base.
\end{description}

% ═══════════════════════════════════════════════════════════
\section{A first session}
{\small\ttfamily
\$ lit conscribe\\
\$ lit inscribe -{}-all -m "first draft"\\
\$ lit account login -{}-email you@ex.com\\
\$ lit push\\
\$ lit show historia\quad\textrm{\itshape (copy a version id)}\\
\$ lit share <version> -{}-expires 7\\
\$ lit publish
}

\vspace{2pt}
\emph{On another machine:} \ci{lit clone <project-id>}, then
\ci{lit pull} to catch up.

\end{multicols}
\end{document}
